
% Flashcard
\pgfkeys{
    /flashcard/.is family, /flashcard,
    default/.style = {},
    color/.store in = \flashcardcolor,
    title/.store in = \flashcardtitle,
    foot/.store in = \flashcardfoot,
}

\NewDocumentCommand{\flashcard}{o}{%
    \pgfkeys{/flashcard, default, #1}%
    \begin{tikzpicture}%
        \pgfmathsetmacro{\w}{(5.715/1pt)}
        \pgfmathsetmacro{\h}{(\w/3*4)}
        \pgfmathsetmacro{\tb}{(\h*0.06)} % top bottom margin
        \pgfmathsetmacro{\lr}{(\w*0.025)}
        \pgfmathsetmacro{\ch}{(\h-\tb*2)}

        % frame
        \begin{scope}
            \clip[overlay,even odd clip,reverseclip] (\lr, \tb) rectangle(\w-\lr, \h-\tb);
            \draw[fill=\flashcardcolor, draw=none, rounded corners=\w*0.04*1cm] (0, 0) rectangle(\w, \h);
        \end{scope}

        % title
        \expandafter\ifstrequal\expandafter{\flashcardcolor}{gray}{
            \def\titlecolor{black}
        }{
            \def\titlecolor{white}
        }
        %% use overlay so that it doesn't affect the bounding box
        \node[align=center, anchor=center, text=\titlecolor, font=\bfseries\sffamily, overlay] at (\w/2, \h - \tb/2) {\flashcardtitle};

        %% proportions: 10%, 30%, 50%, 10%
        % pronounciation
        \node[align=center, anchor=center, minimum height=\ch*0.1] at (\w/2, \ch*0.95+\tb) {
            \begin{tikzpicture}
                \draw[fill=none, draw=black] (0,0) rectangle(\w-\lr*2,\ch*0.1);
                \node[align=center, anchor=center, text=white, overlay] at (\w/2, \ch*0.1/2) {pronounciation};
            \end{tikzpicture}%
        };

        % spelling
        \node[align=center, anchor=center, minimum height=\ch*0.3] at (\w/2, \ch*0.75+\tb) {
            \begin{tikzpicture}
                \draw[fill=none, draw=black] (0,0) rectangle(\w-\lr*2,\ch*0.3);
                % \node[align=center, anchor=center, text=white, overlay] at (\w/2, \ch*0.3/2) {spelling};
                \node(n1)[overlay, tight] at (1, \ch*0.3/2) {\square[color=purple, border=true, scale=0.6]};
                \node(n2)[overlay, tight] at (3, \ch*0.3/2) {\square[color=purple]};
                \draw[-, shorten >=0, shorten <=0] (n1) -- (n2);%
            \end{tikzpicture}%
        };

        % picture
        \node[align=center, anchor=center, minimum height=\ch*0.5] at (\w/2, \ch*0.35+\tb) {
            \begin{tikzpicture}
                \draw[fill=none, draw=black] (0,0) rectangle(\w-\lr*2,\ch*0.5);
                \node[align=center, anchor=center, text=white, overlay] at (\w/2, \ch*0.5/2) {picture};
            \end{tikzpicture}%
        };

        % hint
        \node[align=center, anchor=center, minimum height=\ch*0.1] at (\w/2, \ch*0.05+\tb) {
            \begin{tikzpicture}
                \draw[fill=none, draw=black] (0,0) rectangle(\w-\lr*2,\ch*0.1);
                \node[align=center, anchor=center, text=white, overlay] at (\w/2, \ch*0.1/2) {hint};
            \end{tikzpicture}%
        };

        % footnote
        \node[align=center, anchor=center, text=\titlecolor, overlay] at (\w/2, \tb/2) {\flashcardfoot};

        \useasboundingbox (0,0) rectangle (\w, \h);

    \end{tikzpicture}%
}

\pgfkeys{
    /makeflashcard/.is family, /makeflashcard,
    default/.style = {picture={}, hint={}, tags={}, group={}},
    part/.store in = \fpart,
    % sound
    pronounciation/.store in = \fpronounciation,
    % form
    spelling/.store in = \fspelling,
    inflection/.store in = \inflection,
    % meaning
    translation/.store in = \ftranslation,
    picture/.store in = \fpicture,
    hint/.store in = \fhint,
    tags/.store in = \ftags,
    group/.store in = \fgroups,
}

\NewDocumentCommand{\makeflashcard}{m}{
    \noindent
    \flashcard[color=purple, title=Auge, foot=eye]
    \newpage
}